\section{Order parameter, hidden symmetry and vortices}
\label{sec:order}

\subsection{Frame components and coordinate components}

The order parameter of the superfluid is the expectation value $\Delta^{a}$ of the vector pair
operator (\ref{eq:pairs}) along the extra axis $a\in\{0,5,6,7\}$: four complex numbers. On the
deflating spacetime the index $a$ is a \emph{frame} index: $\Delta^{a}$ is the component along the
unit vector $e_a$. The coordinate components along the deflating coordinates are
$\Delta_t=b_t(x4)\,\Delta^{t}$, and they shrink with the scale factors when the frame components are
constant. All statements below are about frame components.

The energy of the paired state depends on the order parameter through the $8\times8$ gap matrix
$\MM$ on the internal states of a pair. The identity
\begin{equation}
\tfrac18\,\mathrm{Tr}\big[(\MM\MM^\dagger)^{2}\big]
=2\Big(\sum_a|\Delta^{a}|^{2}\Big)^{2}-\Big|\sum_{a,b}\eta_{ab}\,\Delta^{a}\Delta^{b}\Big|^{2},
\label{eq:quartic}
\end{equation}
holds on the deflating spacetime with the frame metric $\eta_{ab}$ of the four axes, to
$2\times10^{-14}$. In coordinate components the same expression contains the metric $g_{tt}$ of the
moment.

\subsection{The ground state: no direction is preferred}

\begin{table}
\caption{\label{tab:states}The potential $V$ of the paired state for six states of the order
parameter, in units of $N(0)\Delta_F^{2}$, where $\Delta_F$ is the gap of a unitary state. Each
state is taken at the magnitude of the order parameter at which its potential is stationary.
$d_-^{2}$ and $d_+^{2}$: the squares of the two fourfold gaps of the state, in units of
$\Delta_F^{2}$. Flat and positive: the numbers of eigenvalues of the matrix of second derivatives of
the potential that vanish and that are positive.}
\begin{ruledtabular}
\begin{tabular}{p{0.37\columnwidth}ccccc}
state & $d_-^{2}$ & $d_+^{2}$ & $V$ & flat & positive \\
\hline
{{ROWS:tbl_states|rows=\SIG{6-11}{0-5}|cols=1,2,3,4,5,6}}
\end{tabular}
\end{ruledtabular}
\end{table}

Table~\ref{tab:states} gives the potential of the paired state for six states of the order
parameter. Every \emph{unitary} state, in which $\Delta^{a}$ is a real unit vector times a phase, has
all eight gaps equal and the lowest potential, whichever direction the vector has. It has four flat
directions, the Goldstone modes (one phase and three rotations of the direction), and four
directions of positive curvature. The \emph{half-paired} states, in which four of the eight gaps
vanish, have half the condensation energy.

One might have expected the deflation to select the direction of the order parameter: the hidden
axis $x0$, which does not deflate, against the three axes that do. \emph{In the sector of fields
that do not depend on the extra coordinates it does not.} The reason is the central observation of
Sec.~\ref{sec:dirac}: in this sector the weighted field obeys the flat equations, and their internal
symmetry rotates the four frame axes into one another. The symmetry of the manifold of ground states
is larger than the symmetry of the spacetime.
%%<1644>

This holds for the model without contact self-interaction, for which Table~\ref{tab:states} is
computed, and for a self-interaction of form Q. \emph{A self-interaction of form C, the definition
of the model, does select a direction.} It breaks the three generators that rotate the hidden axis
into $x5,x6,x7$ (Sec.~\ref{sec:selfsym}), and it enters the pair along the hidden axis and the pairs
along $x5,x6,x7$ with opposite signs, so that their gaps differ (Sec.~\ref{sec:gapformc}). What
selects the direction is the self-interaction, through the distinction that the signature of the
spacetime makes between the space-like hidden axis and the additional times; the deflation still
selects nothing, because it does not act in this sector. The three flat directions of
Table~\ref{tab:states} that rotate the hidden axis into the other three are then no longer flat. The
potential of the order parameter with the contact couplings of form C has not been computed.
%%</1644>

\subsection{The internal symmetry of the sector, and the symmetries of the spacetime}
\label{sec:hidden}

For fields that depend on $x1,x2,x3,x4$ only, the equation of the weighted field contains five
matrices: \SIG{$\alpha^{1},\alpha^{2},\alpha^{3}$, $\beta$ and $i\beta\chi$}{$\alpha^{1},\alpha^{2},\alpha^{3}$,
$\beta$ and $i\beta\chi$}, the last through the coupling of the phase mode. The matrices that commute
with all five generate the \emph{internal symmetry} of the sector. The program computes it again for
the deflating spacetime, item by item (Table~\ref{tab:xsymmetry}), and finds it unchanged:
\begin{itemize}
\item the algebra is $su(2)_A\oplus su(2)_B\oplus u(1)_F\oplus u(1)_E$, of dimension 8;
\item the eight internal states of a fermion form the two doublets $(2,1)+(1,2)$;
\item the four vector pair operators (\ref{eq:pairs}) form the $(2,2)$, with fermion number 2, and
their one invariant quadratic form is the frame metric of the four axes, with the signs
\SIG{$(+,-,-,-)$}{$(+,+,+,+)$ up to an overall sign};
\item a condensate along one axis breaks $SO(4)=[SU(2)_A\times SU(2)_B]/Z_2$ to $SO(3)$.
\end{itemize}
Here $U(1)_E$ is a second phase symmetry of the sector, with a generator $P$ that has the square one
and distinguishes the doublets of $SU(2)_A$ from those of $SU(2)_B$.

Six of the generators are the products $\gamma^{A}\gamma^{B}$ of the six pairs taken from the four
axes $x0,x5,x6,x7$.
%%<1671>
On the flat $7+1$ spacetime they are the rotations of the four extra space-like coordinates, and the
internal symmetry $SO(4)$ is a symmetry of the spacetime. On the deflating spacetime this is no
longer so (Fig.~\ref{fig:isometries}). Of the six, the spacetime keeps
\begin{itemize}
\item none, if the three rates differ;
\item one, the rotation of the two coordinates with the same rate, if two rates are equal;
\item three, the rotations of $x5,x6,x7$ among themselves, if the three rates are equal.
\end{itemize}
The three rotations that mix the hidden axis $x0$ with a deflating coordinate are never symmetries
of the deflating spacetime: $x0$ keeps its size and the other three shrink. The generators that are
lost as symmetries of the spacetime remain exact symmetries of the sector. They are \emph{hidden
symmetries}: symmetries of the equations of a sector that are not transformations of the spacetime.
On the $4+4$ spacetime this is the situation from the start, with or without deflation, because
the rotations that mix the hidden coordinate with the additional times are not unitary there.
%%</1671>
%%<1644>
On the flat $4+4$ spacetime only three of them belong to transformations of the spacetime that the
sector inherits: the products $\gamma^{s}\gamma^{t}$ of two additional times generate the rotations
$SO(3)_T$ of $x5,x6,x7$. The other three, $\gamma^{0}\gamma^{t}$, are Hermitian; the unitary
transformations that they generate in the quantized sector are not the boosts of the planes of $x0$
and $xt$, which are not unitary. The full $SU(2)_A\times SU(2)_B$ was therefore a \emph{hidden
symmetry} already on the flat spacetime: a symmetry of the equations of a sector that is not a
transformation of the spacetime. On the deflating spacetime even $SO(3)_T$ is a
symmetry of the spacetime only if the three rates are equal (Fig.~\ref{fig:isometries}). With two
equal rates one rotation remains, and with three different rates none. The internal symmetry of the
sector is nevertheless exact in all three cases.
%%</1644>

\subsection{Vortices and their zero modes}
\label{sec:vortex}

Because the vacuum value of the Vorton-Scalar breaks the phase symmetry $\UVS$, the field has
line-like defects \cite{NielsenOlesen1973}: along a straight line, say the $x3$ axis, its magnitude
vanishes, and around the line its phase winds by $2\pi$ times an integer. For unit winding,
$\PhiVS$ is $\vVS/\sqrt2$ times a profile that rises from 0 on the line to 1 at a distance of the
order of $1/m_s$, times $e^{i\theta_{12}}$ with the angle $\theta_{12}$ around the line in the plane
of $x1$ and $x2$.

\emph{Energy per unit length.} It is proportional to the square of the four-dimensional vacuum
value and so falls as $e^{-\td}$: in the four-dimensional description vortex lines become lighter as
the deflation proceeds, while their core, of the radius $1/m_s$, keeps its size.

\emph{Fermions on the line.} On the line the fermion mass vanishes, and fermions can be trapped
there \cite{JackiwRossi1981}. In the field of the vortex the Hamiltonian of a fermion is
\begin{equation}
\sum_{j=1}^{3}\alpha^{j}p_j+m\,(\text{profile})\,\beta\,\big(\cos\theta_{12}+i\chi\,\sin\theta_{12}\big),
\end{equation}
with the momenta $p_j$ along $x1,x2,x3$. A state of zero energy at zero momentum along the line is a
constant spinor $w_z$ times a function that decays away from the line. Inserting this form, the
Hamiltonian annihilates the state for every angle if and only if
\begin{equation}
\big(i\alpha^{1}+\beta\big)\,w_z=0,\qquad\big(i\alpha^{2}+i\beta\chi\big)\,w_z=0 .
\label{eq:zeromode}
\end{equation}
The five matrices $\alpha^{1},\alpha^{2},\alpha^{3},\beta,i\beta\chi$ are Hermitian, have the square
one and anticommute with one another. The two conditions state that $w_z$ is an eigenvector, with the
eigenvalue $-1$, of the two commuting matrices $i\beta\alpha^{1}$ and $\alpha^{2}\beta\chi$, and each
joint eigenspace of these two matrices has the dimension $16/4=4$. \emph{A vortex line of unit
winding carries four zero modes of the sixteen-component field.}

For a momentum $p_3$ along the line the energy is the eigenvalue of $\alpha^{3}p_3$. On the four
zero modes $\alpha^{3}$ equals the matrix
\begin{equation}
P_c=i\,\alpha^{1}\alpha^{2}\alpha^{3}\chi,
\end{equation}
the product of the five matrices, which commutes with each of them, is Hermitian and traceless and
has the square one; it is, up to its sign, the generator $P$ of $U(1)_E$. Its trace on the four zero
modes vanishes. \emph{Two zero modes move along the line in one direction at the speed of light and
two in the other}, and the direction of motion is tied to the charge under $U(1)_E$. For winding $-1$
the two directions are exchanged.

Equations (\ref{eq:zeromode}) contain only the matrices of the frame and the fermion mass, neither of
which changes. The four zero modes therefore exist at every moment of the deflation. The count also
holds on flat spacetime. (The familiar count for a four-component field is one zero mode for each
unit of winding; the sixteen components are equivalent to four such fields.)

Massless fermions that run along a vortex line carry a current \cite{Witten1985}; a closed loop of
the line can be held open by that current, and such loops are the \emph{vortons} of the literature
on cosmic strings \cite{DavisShellard1989}. Because the energy per length of the line falls as
$e^{-\td}$ while the fermions on it stay massless, the balance between the tension of a loop and the
current on it shifts in favor of the current as the deflation proceeds. This paper establishes the
ingredients of that statement; it does not compute the equilibrium radius of a loop.

